\chapter{Performance}\label{ch:performance}

A kernel on one G17 core takes at least as long as each of the following limits, and in the measured kernels the
largest of them sets its time.

\begin{enumerate}
\def\labelenumi{\arabic{enumi}.}
\item One simdgroup issues a simple instruction about every 3.1~ns (\cref{sec:scalar-issue}).
\item The tensor unit completes a 16 × 16 × 16 MMA about every 4.96~ns per core at saturation
  (\cref{sec:tensor-issue-rate}).
\item A dependent MMA or a load-to-use chain in one simdgroup costs far more than the issue interval: 74~ns per MMA in
  one chain, 30.5~ns per L1 load and 405–415~ns per DRAM load. Only independent work hides it, from the same simdgroup
  or from other resident simdgroups (\cref{sec:tensor-issue-rate,sec:memory-bandwidth-footprint}).
\item Sustained read bandwidth depends on the footprint: 1.3–4.8~TB/s on-core (a lower bound) and
  269.6–298.5~GB/s from DRAM, depending on the instrument (\cref{sec:memory-bandwidth-footprint}).
\item A hot loop body over 16 KiB loses the instruction cache at low occupancy (\cref{sec:memory-bandwidth-footprint}).
\item Each dependent dispatch costs 3.7–4.2~µs, and each extra compute encoder about 6.5~µs
  (\cref{sec:dispatch-encoder-command-buffer}).
\end{enumerate}

Residency, the number of simdgroups a core holds at once, sets how much independent work other simdgroups can supply
to the third limit (\cref{sec:residency-occupancy-measurements}).

\leadhead{Notation.}

\begin{itemize}
\item \emph{Per-core rate}: summed over every simdgroup on the core. \emph{Per-simdgroup rate}: one instruction
  stream's. The device has 20 cores and four SIMD slots per core (\cref{ch:device-execution-model}).
\item \emph{Occupancy} or \emph{simdgroups per core}: either the number launched divided by 20 or the number resident
  at once; the text says which. A kernel is \emph{saturated} when adding simdgroups per core no longer raises its
  throughput.
\item \emph{Footprint}: the unique bytes a kernel touches.
\item Every absolute time in this chapter is at a warm clock unless stated, timed after a saturating dispatch long
  enough to raise the GPU to its top performance state (\cref{sec:clocks-power-gpu}).
\item MMA opcode names are those of \cref{ch:tensor-unit}, and the tensor MMA used for rate measurements is \op{5106}.
\end{itemize}

Beyond the 0.4–1.3 percent in-process and 9–13 percent cross-process spread of \cref{sec:clocks-performance-states},
one row differed by 18 percent between processes. The retained layer harness, \code{layerbench}, moved 25–44 percent
between its own processes, but ratios timed together in one process reproduce to about 1.5 percent (measured,
MM~\mm{16}, MM~\mm{25.124.2}). Comparisons in this chapter are ratios timed in one process. Every A/B here is also
order-balanced, because under load the second run of a fixed A-then-B pair ran about 0.1~ms per token faster
(measured, MM~\mm{17} rule 43).

\section{Tensor unit}\label{sec:tensor-issue-rate}

At saturation, 16,384 threadgroups × 4 simdgroups × 1,400 issues of the 16-bit MMA, \op{5106}, is 91.75 M issues in
22.44~ms (measured, MM~\mm{8}). That is 4.09 × $10^{9}$ issues per second, or 33.5 TFLOP/s dense fp16/bf16 into fp32 at
8,192 flop per issue and 1.67 TFLOP/s per core. Per SIMD it is 128 MAC per cycle, four times a 32-lane FMA, at an
implied 1.63~GHz. int8, using \op{10384}, runs at about twice the rate, with 13~ns per issue in a chain and 11.50~ms
saturated, which is 8.1 × $10^{9}$ issues per second, 67 TOPS and 256 MAC per cycle per SIMD.

An independent later harness gives 4.96–4.97~ns per MMA per core at 16 simdgroups per core, 8,192 flop in 8.0
cycles at 1.62~GHz, 1,024 flop per cycle per core; int8 2.47~ns. About 5.5 simdgroups per core saturate the unit
(measured, MM~\mm{8}, MM~\mm{0.1}). Percentages of peak here use the later harness's 33.1 TFLOP/s;
the 33.5 derived above is 1.2 percent higher, within harness spread (MM~\mm{8}).

Each core has four tensor slots, one per SIMD (measured, MM~\mm{8}): one, two or four simdgroups each running the
chain take the same time (32.8~µs), 8 take 57.5 and 32 take 224.3 (the full row is in \cref{sec:core-logical-resource}).

\begin{xltabular}{\linewidth}{@{}>{\hsize=1.12\hsize}X>{\hsize=1.25\hsize}X>{\hsize=0.62\hsize}X@{}}
\caption{Cost per MMA by occupancy and by dependence between MMAs.}\label{tab:mma-latency-occupancy}\\
\toprule
Configuration & Cost per MMA & Source \\
\midrule
\endfirsthead
\tablecontinued{3}
\toprule
Configuration & Cost per MMA & Source \\
\midrule
\endhead
dependent chain at saturation & 22.2~ns; eight independent accumulators over the same 1,400 issues take
the same 31.96~µs, so result latency is no longer than the issue
interval & measured, MM~\mm{8} \\*
one simdgroup, 8 independent accumulators & 20.5~ns & measured, MM~\mm{8} \\*
one simdgroup, one dependent chain & 74~ns (3.6 times) & measured, MM~\mm{8} \\*
one simdgroup per core: eight dependent MMAs against four independent chains of two & both 218.0~ns per iteration, the cost
of a single chain & measured, MM~\mm{8} \\*
fit over 321 kernel timings, bf16 class & 62.6~ns inside a dependent chain at 1 simdgroup per core; 5.8~ns at 64 & measured fit, MM~\mm{16} \\*
same fit, fp8 and int8 & 68.2~ns at 1; 7.0~ns at 64 & measured fit, MM~\mm{16} \\*
\op{5107} back-to-back
into the same destination & 34~ns (a write-after-write stall), gone across 8 destinations or at
saturation & measured, MM~\mm{8} \\*
\bottomrule
\end{xltabular}

The same fit charges 1.6~ns per other instruction at one simdgroup and 0.21–0.34~ns at 64. Its six worst outliers are
all independent-MMA ablations, so the large per-MMA figure is the cost of dependence (MM~\mm{16}). The
fp32-operand forms (\cref{sec:instruction-family}) and a transposed A cost nothing extra (measured, MM~\mm{8}).

At 12 or more simdgroups per core interleaved and chained accumulators take the same time (MM~\mm{8}, MM~\mm{11}).

\textbf{Evidence.} MM~\mm{0.1}, MM~\mm{8}, MM~\mm{11}, MM~\mm{16}, MM~\mm{25.144.12}.

\section{Layer kernels}\label{sec:measured-bottlenecks-studied}

An attention chunk loop executes about 46 instructions per MMA: 1,480 instructions for 32 MMAs is 356~ns of the 433~ns
a key block costs per core, while the tensor pipe is busy for 158~ns. Fused attention therefore reaches 35–42 percent
of 33.1 TFLOP/s (11.5–13.8 TFLOP/s) and the FFN 47–57 percent, both bound by issue and the ALU (measured, MM~\mm{8},
MM~\mm{0.12}, MM~\mm{16}). Attention throughput stays at 11.5–13.8 TFLOP/s from an 8 MiB working set to 128 MiB, so
the KV stream is not the limit. With one simdgroup's dependent chain at 5.35~µs per 32-key block, about 247
simdgroups are needed to saturate the GPU (5.35~µs / 433~ns per core × 20 cores).

Simdgroups per threadgroup from 1 to 16 are within noise, provided there are at least 20 threadgroups so every core
has one, and transposed-B flags cost 1.3 percent (measured, MM~\mm{16}). Unrolling K by 2 or 4 does not reduce the
time and adds 98–110 spill stores, and two token tiles per task double the time with 354 spill stores (measured,
MM~\mm{16}). The attention kernel rebuilt on the retained harness costs 1.25–1.29 times the kernel of the
reconnaissance study (\code{docs/g17-tensorops-accelerator-recon.md}, cited as recon below) per key block (6.9
against 5.35~µs for one
simdgroup; 558 against 433~ns per block and core) despite 1,033 instructions and no spills, against the recon
kernel's 1,925 instructions and 42 spill stores (measured, MM~\mm{25.126}).

The runtime K loop is 1.3–1.8 times faster than unrolling at 8 blocks and 5.9–7.2 times at 32 (unrolled kernels
spill 24–923 stores; MM~\mm{16}). A spill-free fused chain is 5–10 percent faster than any cut and 8–20
percent faster than three kernels; at 15 or more spill stores the cuts win. Attention stays fused in all 27 measured
shapes even while spilling 42 stores (cutting P costs 1.03–1.30 times, cutting both 1.14–2.03 times).

Threadgroup memory limits the tile count of a fused gated FFN. The fused form needs \code{nsg*mt*640} bytes for fp8
and int8, or \code{nsg*mt*1,088} for bf16 and fp16, and the 32,768-byte limit caps \code{nsg*mt} at 51 and 30
respectively. At d = 2048 and hidden width 8192, one wave is 20 threadgroups at about 1.8~ms in fp8.

\begin{xltabular}{\linewidth}{@{}>{\hsize=1.15\hsize}Xl>{\hsize=0.85\hsize}Xl@{}}
\caption{GEMM time for M × 128 × 256 half: Apple's \code{matmul2d} against this compiler's K-loop and unrolled forms (measured, interleaved medians, MM~\mm{25.107}).}\label{tab:gemm-vs-matmul2d}\\
\toprule
Shape and grid & Apple & this compiler, K loop & unrolled \\
\midrule
\endfirsthead
\tablecontinued{4}
\toprule
Shape and grid & Apple & this compiler, K loop & unrolled \\
\midrule
\endhead
M512, 8 threadgroups & 106.6~µs & 113.7~µs (1.07 times behind) & 262.7~µs \\*
M1024, 8 threadgroups & 183.6~µs & 244.7~µs (1.33 times) & 376.0~µs \\*
M4096, 256 threadgroups × 1 simdgroup (equal geometry) & 58.7~µs & 62.0~µs (1.06 times) & 78.6~µs \\*
\bottomrule
\end{xltabular}

Apple's own 64 × 4 geometry takes 69.4~µs for M4096, and at 12.8 simdgroups per core the K loop reaches 4,330
GFLOP/s (62.0~µs, interleaved medians, MM~\mm{25.107}) against Apple's 4,570. An earlier occupancy sweep of this kernel
(one 16-row tile row per simdgroup, a 12.5~KB body; round one, which the interleaved medians did not reproduce, so its
absolute figures are lower) goes from 132 GFLOP/s at 8 simdgroups (0.4 per core) to 3,329 at 256 (12.8 per core,
80.6~µs) and 4,497 at 512 (25.6 per core, 124 registers); work grew 32 times and time 1.27 times (measured,
MM~\mm{25.98.1}).
Software pipelining (rotating MMAs ahead of loads) gave 1.08 times at 8 simdgroups (113.8 to 105.6~µs, against a preregistered target of at least
1.2) and 0.99 times at 12.8 per core (measured, MM~\mm{25.97}, MM~\mm{25.107}). In the K-loop
body, one slice per trip waiting on its own loads paid a full round trip per step, about 417~ns against Apple's
245 at 8 threadgroups. A 255-trip one-slice K loop in one simdgroup averages 1.45~µs per trip, with its loads and
four MMAs not overlapping (measured, MM~\mm{25.101}). Issuing two slices' loads before the first MMA (\code{kloop\_unroll=2}) gains 35–38 percent and
comes within 3–4 percent of Apple at 4 and 8 threadgroups (about 337 to 210~ns per step against Apple's 202); an
interleaved control gains 1–4 percent (measured, MM~\mm{25.124.6}).

Unrolled bodies are fetch-bound, at about 518~ns per MMA (512 MMAs in 265~µs) against a 20–74~ns issue interval;
fully unrolled K=256 bodies are 61~KB and 116~KB against Apple's 5,142-byte looping kernel (measured, MM~\mm{25.97},
MM~\mm{25.98}). Grid-split threadgroups run concurrently, since with 512 MMAs per threadgroup time is flat from G = 1 to 8 (\code{t(8)/t(2) = 0.995}), and the null control (2 MMAs) takes
11.5 against 265.3~µs (measured, MM~\mm{25.89}).

For the bf16 layer kernels (measured, MM~\mm{25.43}), 16 simdgroups per threadgroup is fastest for the two- and
three-kernel forms (2.856 and 3.103~ms); 32 is 1.46–1.86 times slower than 16 in every form, against a run spread of
0.007–0.069~ms.

A weight-stationary FFN is 3.7–4.4 times faster than a token-parallel one at 16 and 32 tokens (measured, MM~\mm{16},
MM~\mm{25.126}); on the retained harness fp8 weight-stationary loses from 48 tokens (1.13 times the
token-parallel time) and bf16 from 64 (1.15). The recon model's crossover of about 64 tokens does
not hold for these kernels.

A GELU register epilogue (x * sigmoid(1.702x), six instructions per slot, 4 registers)
takes 15.6~µs at M16 against 52.2~µs for a memory stage (MM~\mm{25.109}). A full-tile one-dispatch attention, whose
softmax is straight-line code, takes 30.5~µs at 2 rows and 104.3~µs at 16 rows (MM~\mm{25.110}).

\textbf{Evidence.} MM~\mm{0.12}, MM~\mm{8}, MM~\mm{16}, MM~\mm{25.43}, MM~\mm{25.89}, MM~\mm{25.101}, MM~\mm{25.97}, MM~\mm{25.98}, MM~\mm{25.98.1}, MM~\mm{25.107}, MM~\mm{25.109},
MM~\mm{25.110}, MM~\mm{25.124.6}, MM~\mm{25.126}.

\section{Scalar issue}\label{sec:scalar-issue}

A simple instruction's per-simdgroup issue floor is about 3.1~ns whatever its class, and fp32 and int32 do not
dual-issue (\cref{tab:scalar-issue-pairs}).

\widetable
\begin{xltabular}{\linewidth}{@{}>{\hsize=1.53\hsize}Xlll>{\hsize=0.74\hsize}X>{\hsize=0.74\hsize}X@{}}
\caption{Scalar issue cost per op in one simdgroup, each class alone and mixed, against the predictions of one shared port and of dual issue: a counted 128-trip loop over eight independent chains, op counts checked from the decoded bytes, each dispatch in its own process (measured, MM~\mm{25.142.6}, MM~\mm{0.15}).}\label{tab:scalar-issue-pairs}\\
\toprule
One simdgroup, ns per op & alone A & alone B & mixed & one shared port predicts & dual issue predicts \\
\midrule
\endfirsthead
\tablecontinued{6}
\toprule
One simdgroup, ns per op & alone A & alone B & mixed & one shared port predicts & dual issue predicts \\
\midrule
\endhead
ffma + iadd & 4.18 & 3.09 & \textbf{3.61} & 3.64 & 2.09 \\*
fadd + iadd & 3.11 & 3.09 & 3.11 & 3.10 & 1.55 \\*
fmul + shl\_imm & 3.16 & 6.10 & 3.17 & 4.63 & 3.05 \\*
fadd + fmul (control, one class) & 3.11 & 3.16 & 3.17 & 3.14 & 1.58 \\*
\bottomrule
\end{xltabular}

The ffma + iadd pair, which separates the two models, lands on the shared-port value, and the
full grid gives fadd 1,050, iadd 1,238 and mixed 1,271~µs at B = 128. fp32 ffma at 4.18 shows that eight
chains do not fully hide its latency. \code{shl\_imm} costs 6.1 alone and 3.17 mixed with fmul, so its extra cost is a
stall that other instructions can fill (the narrow register file's constraint, \cref{ch:register-file}). fp32 add
throughput is about 4.1 × $10^{12}$ per second machine-wide, about 2.1 × $10^{11}$ per core, consistent with 128 fp32
lanes per core at about 1.6~GHz (the clock was not measured by this probe).

Eight independent chains reach the floor (MM~\mm{17} rule 38). Apple's wait after every
special-register read costs nothing measurable (\cref{sec:read-latency-whether}, MM~\mm{25.121}).

\code{fdiv} costs 9 more instructions than a single-instruction operation such as \op{1272} (measured, MM~\mm{6}); the
other ALU cases that affect cost are in \cref{ch:scalar-arithmetic}.

\textbf{Evidence.} MM~\mm{0.15}, MM~\mm{6}, MM~\mm{17}, MM~\mm{25.121}, MM~\mm{25.142.6}, MM~\mm{25.144.12}.

\section{Register cost}\label{sec:registers-spills-occupancy}

Against a 14-accumulator kernel, 15 accumulators (32 spill stores) cost 1.31 times, 16 cost 1.47, 18 cost 2.17 and
24 cost 2.30 at 1–4 simdgroups per core, but only 1.00–1.10 times at 64 (measured, MM~\mm{11}; register-file rules in
\cref{sec:capacity-budgets-spills}). Removing spills did not always help; with few stores the streaming order ran
at 0.47–0.79 of the flat order's speed in several shapes.
Spill counts here include both spill forms (\cref{sec:capacity-budgets-spills}; MM~\mm{11}, MM~\mm{13}, MM~\mm{19}).

For this compiler's one-tile GEMM body (measured, MM~\mm{25.111}), declaring 44 or 124 registers (a one-byte metadata
change) gives equal time within 1.2 percent at 25.6 simdgroups per core (41.6 against
42.1~µs) and 0.7 percent at 51.2 per core (69.9 against 69.4~µs). Equal time past 32 simdgroups per core does not by
itself show residency past 32. This compiler's GEMM tiles peak at R42 (16 × 16 × 64), R82 (32 × 32 × 64) and R122
(64 × 64 × 64), M512 with GELU reaches R125, and the declared counts are 44, 60, 92 and 124 at N16, 32, 64 and 128
(this compiler, MM~\mm{25.98.1}, MM~\mm{25.109.1}).

In the decode qmv kernel at 4 rows per simdgroup, the same loop at 81 registers instead of 114 is 4.1~µs faster
(43.2 against 47.3~µs; measured, MM~\mm{25.144.4}). The B~8 batched qkv kernel uses 109 (q4)
and 108 (q8) registers and holds occupancy at 25–27 percent where its threadgroup memory (64–128 bytes) and
threads per group would allow 96 simdgroups per core, so the register file sets the cap. Halving its accumulators moved the
count only from 109 to 107 and did not lift occupancy, because its registers are dominated by weight words, dequantized values,
field temporaries and addresses. Instruments' "Register Pressure Influence" counter read 0 for every kernel recorded,
including these.

\textbf{Evidence.} MM~\mm{10}, MM~\mm{11}, MM~\mm{13}, MM~\mm{19}, MM~\mm{25.98.1}, MM~\mm{25.109.1}, MM~\mm{25.111}, MM~\mm{25.115}, MM~\mm{25.144.4}.

\section{Measured residency}\label{sec:residency-occupancy-measurements}

Metal's public API has no occupancy counter; it offers only \code{GPUTimestamp} (measured, MM~\mm{25.48}; \cref{sec:dispatch-encoder-command-buffer}).
\code{maxTotalThreadsPerThreadgroup} reads 1,024 for every pipeline whatever its register and threadgroup-memory
load (\cref{ch:device-execution-model}), and \code{staticThreadgroupMemoryLength} tracks the request (measured,
MM~\mm{25.47}, MM~\mm{11}). Residency was therefore measured
directly with the probes in \cref{tab:residency-counts} (the residency rules are in \cref{ch:device-execution-model}),
and later read through Instruments counters (\cref{sec:dispatch-encoder-command-buffer}).

Time is flat from 20 to 116 registers per lane. Turner's M1/M2 benchmarks, reported in recon section 167, put the ALU
utilization ceiling at about 24 simdgroups per core, and above it occupancy can fall with no change in time, so the
flat curve does not measure occupancy. A direct atomic counter (arrival add, peak max, fixed spin, exit decrement) reads occupancy falling from 44.6
to 32.1 simdgroups per core over 22 to 126 registers (MM~\mm{19}, MM~\mm{20}, MM~\mm{25.48}, MM~\mm{25.49}); the
earlier slot-limited reading of recon section 167 is in \cref{app:a}.

\begin{xltabular}{\linewidth}{@{}X>{\hsize=0.65\hsize}X>{\hsize=1.77\hsize}X>{\hsize=0.61\hsize}X@{}}
\caption{Resident simdgroups measured by each probe.}\label{tab:residency-counts}\\
\toprule
Kernel and probe & Conditions & Resident simdgroups & Source \\
\midrule
\endfirsthead
\tablecontinued{4}
\toprule
Kernel and probe & Conditions & Resident simdgroups & Source \\
\midrule
\endhead
Apple-compiled direct atomic counter (recon 170) & 22 to 126 registers & 44.6 falling to 32.1 per core (1.39 times); about 40 at 44 registers (44.5 at
30, 39.6 at 46); 32.0 at 98 & measured, MM~\mm{19}, MM~\mm{25.49}, MM~\mm{25.111.2} \\*
Recon 134 kernel, 1,024-thread threadgroups & 40 threadgroups at 126 registers & peak exactly 640 resident, 32 per core; resident falls from about
890–930 at 22 registers to 530–660 at 126 (25–30 percent, no step),
worth 1.00 to 0.95 in throughput capacity & measured, MM~\mm{11} \\*
This compiler's residency counter, straight-line, 12,800
instructions, one simdgroup per threadgroup & warm clock, launched past capacity (S = 4,096), 3 rounds & 95.0 per core at 19 registers used; 71–74 at 33; 67–69 at 49; 66–69 at
65; 65–66 at 81; 64–66 at 97; 68–69 at 113; 64–68 at 125 & measured, MM~\mm{25.121} \\*
Same probe, declared against used & S = 1,024 and 2,048 & declaring 125 while using 19 peaks with using 19 (1,814–1,896 at S =
2,048); warm, using 125 holds 1,024 of 1,024 in 5 of 5 & measured, MM~\mm{25.115} \\*
Threadgroup-memory handle, 128-thread threadgroups, 4,096 threadgroups & 9 to 94 registers & at least 8 simdgroups per core at 94 registers (two threadgroups still
co-resident) & measured, MM~\mm{25.49} \\*
Forward-progress probe, single-simdgroup threadgroups & spinning residents & up to 1,200 threadgroups progress concurrently; first resident wave at
2,000 launched about 1,690 (about 84 per core at that register count);
the 1,024-thread runs, and the restriction they motivate (do not depend
on preemption, and size a waiting grid to residency), are in
\cref{sec:forward-progress-in-order} & measured, MM~\mm{25.141.1} \\*
Instruments, B~8 batched qkv decode kernel & 109 registers & occupancy 25–27 percent of a 96-simdgroup-per-core budget & measured counters, MM~\mm{25.144.4} \\*
\bottomrule
\end{xltabular}

The Apple-compiled counter (32.1 at 126) and the recon-134 kernel (32 per core at 126) agree on a floor near 32 for
Apple-compiled kernels at high register counts; MM~\mm{11}'s gradual 25–30 percent fall and the counter's 1.39-fold
fall describe the same shape on different kernels. This compiler's counter reads about twice as many (64–69 per core
from 49 to 125 registers) on a different kernel, straight-line code from its own allocator. Why the two kernels differ
is open (MM~\mm{25.121}). \code{agxforge/g17/tensorsched.py} uses a separate floor for each kernel family: 64 per
core from 49 to 125 registers for this compiler's kernels, and 32.1 at 126 for Apple's.

Residency follows the registers a program uses, and the count it declares has no effect (rule 1 of
\cref{sec:residency-occupancy-rules}; MM~\mm{25.115}). A fixed-pool model, one pool divided by demand, predicts 7.8
simdgroups at 126 registers where 32.1 are measured. The physical capacity and allocation granularity are unknown (\cref{sec:residency-occupancy-rules}; MM~\mm{25.49}, MM~\mm{20}).

Threadgroup memory also acts as a residency handle (measured, MM~\mm{25.49}). The 57,344-byte per-core budget (rule 4 of
\cref{sec:residency-occupancy-rules}; the 32,768-byte per-threadgroup limit is in \cref{ch:memory}) shows as a step: 28,672 bytes per threadgroup runs 0.00304 s, 28,688 bytes 0.00382 s (+25 percent for 16
bytes). The step persists at 94 registers (ratio 1.047), and a 24-fma control at 10 registers gives 1.041, so
registers do not set its size; it tracks arithmetic intensity.

\subsection*{Per-core simdgroup staircase}

Time per iteration (8 legacy issues) rises in steps with simdgroups per core (\cref{tab:simdgroup-staircase}).

\widetable
\begin{longtable}{@{}rr@{}}
\caption{Median time per iteration of the legacy MMA loop against simdgroups per core. The probe is a retained
Apple-compiled legacy \code{simdgroup\_multiply\_accumulate} loop with 8 independent accumulators and 2,048 trips;
the grid is 20 × simdgroups per core with threadgroups of 32; the clock is warm, with a dispatch of at least 4.7~ms
before every timed one and a fixed 4-simdgroup-per-core reference dispatch after each (measured,
MM~\mm{25.122}).}\label{tab:simdgroup-staircase}\\
\toprule
Simdgroups per core & ns per iteration \\
\midrule
\endfirsthead
\tablecontinued{2}
\toprule
Simdgroups per core & ns per iteration \\
\midrule
\endhead
1–8 & 24.5–24.6 \\*
9–12 & 30.7–31.2 \\*
13–17 & 54.6 \\*
18–21 & 60.2–60.9 \\*
22–30 & 83.6–84.3 \\*
31–34 & 85–90.6 \\*
35–42 & 113.4–114.5 \\*
43–45 & 119.6–120.7 \\*
46–48 & 143.0–143.6 \\
\bottomrule
\end{longtable}

The rungs sit at 24.6 + 29.6k ns (k = 0..4: 24.6, 54.0–54.6, 83.7, 113.3, 142.9), each with a +6~ns half-step about 4
simdgroups before the next full rung (30.6, 60.2, 89.9, 119.5), and a full rung spans about 11 simdgroups per core.
The first two boundaries fall exactly per core, at 161 and 241 threadgroups (20 × 8 + 1, 20 × 12 + 1). Later
boundaries are windows (358–368, 442–460, 613–641, 675–686, 840–855, 891–923 threadgroups) inside which each dispatch
of the same grid lands on one of two adjacent rungs. The reference reads 24.54~ns either way, which rules out the clock,
and outside the windows one grid repeats within 0.33 percent.

A linear model of total work does not fit, since the steps of 6–24~ns stand well above a 0.5 percent noise floor.
Neither does a threadgroup granule, since threadgroup sizes 32–128 are equal within 1.5 percent at equal work. A
residency or wave granule fits the timings. The rung follows the simdgroup count, and the number of accumulators does not move it: with 2, 4, 8 or 16 accumulators
per simdgroup the boundaries sit at the same simdgroups per core within one, and ns per MMA at 1–8 simdgroups per core
is 11.87, 5.99, 3.06 and 1.61, a constant about 24~ns per iteration, so up to 16 independent chains hide behind one
dependent-chain latency. Writing time = (one simdgroup's iteration time) × f(simdgroups per core), f is about 1 up
to 8, 1.25 at 9–12, 2.2 at 13–17 and 2.5 at 18–21. Of the two granules that
fit, a residency cap is ruled out, since the residency counter holds 64–95 per core (MM~\mm{25.121}). The cause
inside the core is unknown.

A "bimodal legacy rate" that flipped from dispatch to dispatch (open item U3) was first read as a
two-state per-issue rate with a constant 9/8 ratio (recon 158), then as a 79.4~ns quantum (recon 164, in fact the
slope of a line). Both readings came from measurement errors. The occupancy axis was off by a factor of 32 and the
grid was oversubscribed (MM~\mm{25.85}, MM~\mm{25.86}), and runs shorter than about 1.5~ms ran at the idle clock (a v2 sweep with a 0.5~ms
warm-up saw medians of 29.9 and 49.6~ns in the two halves of one run). On the corrected axis at a warm clock the
bimodality is confined to the boundary windows above (MM~\mm{25.122}, MM~\mm{19}).

A legacy MMA chained into one accumulator measures latency, and independent chains scale the rate until their
accumulators outgrow the registers. One chain gives 129.5 GMAC/s, 2 give 259.1, 4 give 518.1, then 8 and 16 fall to 395.6 and 199.0 (measured,
MM~\mm{25.56}).

\textbf{Evidence.} MM~\mm{11}, MM~\mm{19}, MM~\mm{20}, MM~\mm{25.47}, MM~\mm{25.48}, MM~\mm{25.49}, MM~\mm{25.56}, MM~\mm{25.111.2}, MM~\mm{25.115}, MM~\mm{25.121},
MM~\mm{25.122}, MM~\mm{25.141.1}, MM~\mm{25.144.4}.

\section{Memory and instruction fetch}\label{sec:memory-bandwidth-footprint}

Sustained read bandwidth by footprint (measured, MM~\mm{7}, streaming microbenchmark) is 1.3–4.8 TB/s up to 5 MiB
(on-core, reported as a lower bound), 520–528~GB/s at 10–20 MiB (the SLC plateau) and 281–294~GB/s from 40 MiB
(DRAM), with the step between 20 and 40 MiB. The figures do not change with threadgroup shape at 1, 8 and 32
simdgroups.

\begin{xltabular}{\linewidth}{@{}>{\hsize=1.20\hsize}X>{\hsize=1.22\hsize}X>{\hsize=0.58\hsize}X@{}}
\caption{DRAM read bandwidth by measuring instrument.}\label{tab:dram-bandwidth}\\
\toprule
Instrument & Read & Source \\
\midrule
\endfirsthead
\tablecontinued{3}
\toprule
Instrument & Read & Source \\
\midrule
\endhead
streaming microbenchmark by footprint, 40 MiB and above & 281–294~GB/s (MM~\mm{24.2} row M7: 283–294) & MM~\mm{7}, MM~\mm{20} \\*
1,024-thread threadgroups, 256~MB cold streams & 298.5~GB/s (peak; supersedes the 272~GB/s used before MM~\mm{25.141.13}) & MM~\mm{0.15}, MM~\mm{25.141.13} \\*
grid-stride float4, 1,048,576 threads, 256 MiB and 1 GiB, two rounds & 269.6–275.7~GB/s; write 241.0–252.9; copy 243.5–255.6 (read plus write
counted); 64 MiB reads 218–264 & MM~\mm{25.121} \\*
static performance model floor & 280~GB/s (a model constant) & MM~\mm{25.144.12} \\*
\bottomrule
\end{xltabular}

Roofline figures here use 298.5~GB/s as the peak (\cref{tab:dram-bandwidth}). The grid-stride and streaming instruments agree within 6
percent (MM~\mm{25.121}).

\begin{xltabular}{\linewidth}{@{}XX@{}}
\caption{Load-to-use latency of a dependent load by footprint: 32 lanes each chasing a random cycle, warm clock,
least-squares slope over five chain lengths (measured, MM~\mm{25.121}).}\label{tab:load-to-use-latency}\\
\toprule
Footprint & ns per dependent load \\
\midrule
\endfirsthead
\tablecontinued{2}
\toprule
Footprint & ns per dependent load \\
\midrule
\endhead
16 KiB (L1) & 30.5 (about 49 cycles) \\*
4 MiB & 195 (192–238) \\*
1 GiB, 16k–64k hops (DRAM) & 405–415 \\*
\bottomrule
\end{xltabular}

The static model uses 30 / 96 / 195 / 313 / 422~ns at 16 KiB / 256 KiB / 4 MiB / 64 MiB / 256 MiB (MM~\mm{25.144.12}); the
256 KiB and 64–256 MiB rows do not fit one line in the probe and are model inputs. Issuing a load costs 9.3~ns in an
unwaited chain (measured, MM~\mm{25.97}).

Decode attention keeps K and V resident in the last-level cache between dispatches up to about 4.2~MB (1,024 keys),
at 18–20~ns per key; past it the cost is about 29~ns per key and 140~GB/s, under half the DRAM peak
(MM~\mm{25.181}). Spare threadgroups warming the cache save up to 10~µs per small-to-big kernel pair (7 percent) when
the warm set fits the small kernel's duration (4~MB with 256 extra threadgroups: 143.2 to 133.3~µs). Decoy loads gave
no saving, and 8~MB lengthens the small kernel (MM~\mm{25.141.5}).

\subsection*{Instruction fetch in hot loops}

\begin{xltabular}{\linewidth}{@{}>{\hsize=0.81\hsize}X>{\hsize=1.66\hsize}X>{\hsize=0.53\hsize}X@{}}
\caption{Instruction-fetch cost against loop size, from two sweeps and two follow-up measurements.}\label{tab:instruction-fetch}\\
\toprule
Measurement & Result & Source \\
\midrule
\endfirsthead
\tablecontinued{3}
\toprule
Measurement & Result & Source \\
\midrule
\endhead
Per-instruction time against loop code size, coarse sweep & grows sixfold from 1~KB to 26~KB at one simdgroup per core (0.411 to
2.49~ns), 1.85-fold at 64, then flattens & measured, MM~\mm{7} \\*
Location of the step, seventeen-point sweep at fixed instruction mix & the variable is the loop body: 15,971 bytes fast, 16,595
partial, 17,247 slow, bracketing 16 KiB (1,363 bytes of the
kernel are prologue) & measured, MM~\mm{19} \\*
Cost of crossing it & in recon 159's microbenchmark, a loop body over 16 KiB costs 2.81 times on one simdgroup and
1.15 times at 64 simdgroups per core; cold straight-line code costs 1.91~ns per
byte for one simdgroup & measured, MM~\mm{25.144.12} \\*
At high occupancy, in real kernels & at 26–64 simdgroups of the same code per core fetch is shared and
crossing cost nothing beyond the instruction count: register-O attention bodies of 21.6–22.5
KB folded to 13.1~KB ran within 1.3 percent of the instruction-count
ratio (1,169 against 739~µs = 1.58; 1,829 against 1,141 instructions = 1.60) & measured, MM~\mm{25.144.12} \\*
\bottomrule
\end{xltabular}

A hot loop body should stay under 16 KiB (the fast point is 15,971 bytes),
in place of the earlier guideline of "about 26 KiB, roughly 2,800 instructions" (MM~\mm{0.1}, MM~\mm{7}). Between 1
and 26 simdgroups per core fetch has not been measured (\cref{tab:instruction-fetch}). Neither fetch term improved the
static model (\cref{sec:cost-models-their}).

\textbf{Evidence.} MM~\mm{0.15}, MM~\mm{7}, MM~\mm{19}, MM~\mm{20}, MM~\mm{25.97}, MM~\mm{25.121}, MM~\mm{25.141.5}, MM~\mm{25.141.13}, MM~\mm{25.144.12}, MM~\mm{25.181}.

\section{Dispatch}\label{sec:dispatch-encoder-command-buffer}

\begin{xltabular}{\linewidth}{@{}>{\hsize=0.84\hsize}X>{\hsize=1.33\hsize}X>{\hsize=0.83\hsize}X@{}}
\caption{Measured dispatch, encoder and command-buffer costs through Metal. Submission costs below Metal, including
the resident runtime's warm Submit-to-callback split, are in \cref{sec:submit-record-call,sec:completion-callbacks-their}.}\label{tab:dispatch-costs}\\
\toprule
Item & Cost & Conditions \\
\midrule
\endfirsthead
\tablecontinued{3}
\toprule
Item & Cost & Conditions \\
\midrule
\endhead
dependent tiny dispatch & 3.7~µs each in a serial encoder; 4.2~µs with a concurrent encoder plus
scoped or resource barriers (MLX's pattern) & chain of one-thread dispatches each reading the previous word (MM
\mm{25.141.10}, \mm{25.141.13}) \\*
extra compute encoder in one command buffer & about 6.5~µs each & 172 encoders per token: 9.275 against 8.163~ms as one encoder (MM
\mm{25.142.1}) \\*
concurrent encoder with scoped buffer barriers & 0.03–0.07~ms per token slower than serial & pipelined q8 token, 196 dispatches (MM~\mm{25.142.1}) \\*
tracked buffers in one serial encoder against untracked buffers ordered
by \code{MTLSharedEvent} & tracked faster by 0.2–0.3~ms per token & ABBA, both precisions (MM~\mm{25.142.1}) \\*
indirect command buffer & direct against indirect 4.82 against 5.83~µs per dispatch (64 × 1
threadgroups), 4.66 against 5.81 (64 × 32); no gain for large dispatches
(143.98 against 145.59~µs; 131.80 against 128.41 at 32~MB) & MM~\mm{25.141.7} \\*
pipeline change & not measurable: one uber pipeline for FFN and w2 behind skip branches
101.9~µs per dispatch shared, 101.8 separate, 99.9 same-pipeline floor & MM~\mm{25.141.2} \\*
grid barrier instead of dispatch boundaries & 133.1 against 120.8~µs per 32~MB phase (fenced, capped barrier among
resident threadgroups; no barrier reached its cap in these runs, 
\cref{sec:forward-progress-in-order}) & MM~\mm{25.141.1} \\*
\bottomrule
\end{xltabular}

Removing a pipeline change saved 0.6–2.5~µs, against the 6.54~µs a constant-drain fusion model predicted, so a
pipeline change is not counted as a cost (MM~\mm{17} rule 39); a fusion saves about the work it removes (measured,
MM~\mm{25.142.4}). An in-chain layer A/B predicted the graph at 0.75–1.07 times for changes of 8–22~µs per layer, and
small savings reached the graph only on the path every token waits for. Changes are therefore judged in the whole graph
(MM~\mm{17} rule 41).

A prefill GEMM grid with row groups in the high bits of the threadgroup id ranked differently alone and in the graph. It
was 9 percent faster alone with weights resident and 15 percent slower in the graph, because row groups reading one
column block were \code{grid\_n} threadgroups apart and each re-streamed it from DRAM (measured, MM~\mm{25.157}).

\subsection*{Counters by instrument}

\begin{xltabular}{\linewidth}{@{}>{\hsize=0.81\hsize}X>{\hsize=1.63\hsize}X>{\hsize=0.56\hsize}X@{}}
\caption{Performance counters available through each instrument.}\label{tab:counter-instruments}\\
\toprule
Instrument & What it gives & Source \\
\midrule
\endfirsthead
\tablecontinued{3}
\toprule
Instrument & What it gives & Source \\
\midrule
\endhead
Metal's public counter API & one counter set, \code{timestamp} (\code{GPUTimestamp}), sampled at
stage (encoder) boundaries only: \code{supportsCounterSampling} stage 1,
dispatch 0; sampling at a dispatch boundary asserts & measured, MM~\mm{25.48}, MM~\mm{25.141.6}, MM~\mm{0.15} \\*
Metal System Trace & command buffers and encoders only; the only GPU counter is "RT Unit
Active"; it records the GPU performance state (Minimum, Medium, Maximum)
per interval & measured, MM~\mm{0.15}, MM~\mm{25.142.1} \\*
Instruments \code{xctrace record} with a GUI-saved "Performance
Limiters" template (\code{counterscounterprofile = 13},
\code{g17-limiters.tracetemplate}) & 89 counters, including Kernel Occupancy, Compute SIMD Groups Inflight,
Instruction Throughput Limiter and Utilization, ALU / F32 / F16 /
Integer limiters, L1 Cache Limiter, Buffer L1 Read/Write Bandwidth, GPU
Read/Write Bandwidth, Last Level Cache Bandwidth and Limiter, Memory
Request Stalls Influence, Register Pressure Influence, Occupancy Manager
Target, Neural Accelerator Limiter & measured, MM~\mm{25.144.4} \\*
\bottomrule
\end{xltabular}

The private \code{GPURawCounter.framework} route fails: the sampler-source query (selector 261) returns
kIOReturnNotPrivileged, and root does not help; it needs the entitlement \code{com.apple.private.agx.performance-spi}, and
Apple's \code{MTLReplayer}, which carries it, is killed on launch by AMFI (measured, MM~\mm{25.144.4}).

Per-core Instruments counters are medians over busy samples (Kernel Occupancy at least half its own 90th percentile;
MM~\mm{25.144.4}); bandwidths are 90th percentiles from a separate sampler. Absolute "GPU read
bandwidth" is not a DRAM figure when a looped weight stays LLC-resident (MLX's about 500~GB/s exceeds the 298.5~GB/s
DRAM peak).

No hardware clock is readable inside a kernel: none of the 67 special-register indices of
\op{14059} is a clock, \code{llvm.readcyclecounter} crashes the compiler service, and a GPU-side
ticker slows five-fold under load. The working clock is the CPU's \code{mach\_absolute\_time} (41.667~ns per tick) written
into 512 lines of an untracked write-combined buffer, read by a one-thread stamp dispatch (about 1~µs each); the
stamped span matches the command buffer to 0.14–0.5 percent (measured, MM~\mm{25.141.13}). One encoder per dispatch is not
a faithful timeline (a layer read 276~µs that way against 330 in one encoder; measured, MM~\mm{0.15}; MM~\mm{17}
rule 42).

\textbf{Evidence.} MM~\mm{0.15}, MM~\mm{17}, MM~\mm{25.48}, MM~\mm{25.141.1}, MM~\mm{25.141.2}, MM~\mm{25.141.6}, MM~\mm{25.141.7}, MM~\mm{25.141.13},
MM~\mm{25.142.1}, MM~\mm{25.142.4}, MM~\mm{25.144.4}, MM~\mm{25.157}.

\section{GPU clock and power}\label{sec:clocks-power-gpu}

A cold one-threadgroup pointer chase was linear in chain length (206, 405, 802 and 1,640~µs at 4 MiB) and then fell
for longer chains (measured, MM~\mm{25.121}, MM~\mm{25.122}). The top clock, the governor, the 1.5~ms idle-clock
threshold and the warm-up rule are in \cref{sec:clocks-performance-states}. Short command buffers also drop the
performance state, and mlx-lm's seven per token spent 27–78 percent of decode in "Medium" (MM~\mm{0.15}). A clock
reference does not detect contention for cores. Under another workload (Device Utilization 31 percent) the same FFN
projection timed 185–194~µs against 145–155 idle, while the reference dispatch read 50.3–50.6~µs in both (measured,
MM~\mm{25.124.6}).

The rows of \cref{tab:gpu-rail-power} from MM~\mm{25.56} and from MM~\mm{25.60} used different instruments and
engines, so no ratio is taken across the two sets.

\begin{xltabular}{\linewidth}{@{}>{\hsize=1.48\hsize}X>{\hsize=0.91\hsize}X>{\hsize=0.91\hsize}X>{\hsize=0.70\hsize}Xl@{}}
\caption{GPU-rail power, work rate and efficiency for each arm.}\label{tab:gpu-rail-power}\\
\toprule
Arm & GPU rail & Work & Efficiency & Source \\
\midrule
\endfirsthead
\tablecontinued{5}
\toprule
Arm & GPU rail & Work & Efficiency & Source \\
\midrule
\endhead
idle & 32.3 mW & – & – & MM~\mm{25.56} \\*
fp32 fma & 4,948 mW & 12.4 GMAC/s & 2.5 GMAC/J & MM~\mm{25.56} \\*
legacy engine, \op{838}, four
independent accumulators & 9,973 mW & 184.4 GMAC/s & 18.5 GMAC/J & MM~\mm{25.56} \\*
legacy engine, \op{838}, one dependent
chain & 19,001 mW & 129.6 GMAC/s & 6.8 GMAC/J & MM~\mm{25.56} \\*
idle (second instrument) & 31.6 mW & – & – & MM~\mm{25.60} \\*
fp32 FMA, serially dependent chain & 27,917.5 mW & 0.071 TFLOP/s (latency-bound) & – & MM~\mm{25.60} \\*
TensorOps \code{matmul2d} 32 × 32 × 64 half to float & 62,768.4 mW (reproduced at 62,518.0, 0.40 percent) & 4.00 TFLOP/s, 12 percent of 33.5 & 2.06 uJ per 32 × 32 × 64 MMA & MM~\mm{25.60} \\*
\bottomrule
\end{xltabular}

The MM~\mm{25.56} rows used 5 s windows, three trials and a grid of $2^{20}$, with the GPU at 1,620~MHz and 100
percent active residency in every working arm. The dependent chain draws 1.9 times the power of the four-chain kernel
for less work (trials within 40 mW). The fp32 fma arm draws 4,948~mW; the repository README gives 5,607~mW for a
kernel computing and 1,681~mW for one stalled on memory.

The MM~\mm{25.60} rows used filled 6 s windows, with liveness ratios (runtime against an eight-fold loop increase) of
8.05 for the tensor kernel and 7.82 for the ALU. The 2.25-fold ratio between them (62,768 against 27,917~mW) compares
a latency-bound ALU chain with a tensor kernel at 12 percent of peak; neither row is a peak. The first tensor attempt in MM~\mm{25.60}
stored its result under \code{if (lane == 1024)}, the compiler deleted the loop, and it read 9.5~W; dispatch time did
not move with the loop count (256 and 32,768 iterations both 0.010~ms), which showed that the loop was gone.

Killing the host of a running kernel left the kernel running and the GPU at 1,620~MHz, 5,588 mW and 0 percent idle,
with every process reporting 0 GPU ms/s (MM~\mm{17} rule 35); recovery was a reboot (\cref{sec:faults-hangs-recovery}).

\textbf{Evidence.} MM~\mm{0.15}, MM~\mm{8}, MM~\mm{17}, MM~\mm{20}, MM~\mm{25.56}, MM~\mm{25.60}, MM~\mm{25.121}, MM~\mm{25.122}, MM~\mm{25.124.6}.

\section{Cost models}\label{sec:cost-models-their}

The scheduler model, \code{tensorsched}, prices uniform work in measured kernel families and bounds non-uniform work
from below. The static model, \code{g17perfmodel}, prices any compiled program from its bytes.

\subsection*{Scheduler model (\code{agxforge/g17/tensorsched.py})}

The model (MM~\mm{16}, MM~\mm{25.124}, and the module docstring) is:

\begin{lstlisting}
T_latency    = max over simdgroups of (steps of that simdgroup) x L_step
T_throughput = (total steps) x C_step / cores
T_memory     = unique bytes / BW(footprint)
time         = max(T_latency, T_throughput, T_memory)
\end{lstlisting}

This is the standard makespan lower bound, tight only when every simdgroup does the same number of steps.

On uniform work (measured, MM~\mm{16}, MM~\mm{25.124}) the attention median error is 2.7 percent (20 of 23 rows within
10 percent; worst −16.5 percent at 128 × 2,048 keys fp8). Weight-stationary stage 1 is stated at 5.1 percent and gives
3.3 percent on the retained rows (worst −6.0). The token-parallel FFN is stated at 3.3 percent over 16 rows, of which 7
come from an unretained log; the 9 retained rows give 4.7 percent (worst +11.9, bf16 at 10 threadgroups).

A causal attention gives query tiles triangular trip counts, and there the model over-predicts the diagonal skip's
benefit by 20.9 percent at 256 tasks and 37.1 percent at 1,024; refitting both constants on the same kernels' uniform
runs leaves the error (measured, MM~\mm{16}). For unequal work its predicted speedup is a ceiling and its time a lower
bound. The uniform two-regime shape is also too sharp (time rises 1.46 times for 8 times the tasks, then 3.14 times
for the next 4 times; the best fit leaves ±12 percent).

A scheduling term, \code{T\_schedule = list\_makespan\_ns}, narrows the gap (model and measured, MM~\mm{25.124.1}).
Tasks are dealt in threadgroup index order round-robin to 20 cores, then one per freed slot, and a core with n
resident tasks advances each one step per \code{max(L\_step, n x C\_step)}. On the causal points (e4m3, head width
128, \code{L\_step} 2,868~ns, \code{C\_step} 327~ns) speedup optimism falls from 20.9 to 7.5 percent at 256 tasks and
from 37.1 to 26.5 percent at 1,024. The term is still a lower bound, and \code{speedup} refuses to form a ratio from it.

The model's constants (MM~\mm{25.124}) are the attention step pairs 5.35~µs / 433~ns (fp8, int8) and 5.0~µs / 381~ns
(bf16), and a bandwidth ladder of 3 TB/s to 5 MiB, 0.52 TB/s to 20 MiB and 0.28 TB/s from 40 MiB. The chain fit is
62.6 / 5.8~ns (bf16) and 68.2 / 7.0~ns (fp8, int8), and \code{chain\_ns} refuses any occupancy but 1 and 64 per core.
The residency floors are 64 per core from 49 to 125 registers (this compiler's kernels) and 32.1 at 126 (Apple's).
Every prediction carries checks (kernel family, dtype, head width or layer shape, registers 19–126 with the residency floor holding
the needed simdgroups, footprint on a measured plateau, token range, uniform work); a failed check yields no number,
and \code{choose} returns the known-safe schedule with \code{fallback=True}.

On the retained harness the recon constants chose the wrong FFN
schedule at fp8 48 and bf16 64 tokens, while the \code{layerbench} constants pick the measured-fastest
at all 11 points (blocked weight-stationary to token-parallel between 32 and 48 e4m3 tokens and between 48 and 64
bf16; one-pass weight-stationary wins at 48 e4m3, 1.119 against 1.254~ms). Head width scales by the MMA count
(\code{dh/128} for multiples of 32 from 96 to 256, 12 ratios at −6.4 to +2.3 percent; 144 costs 1.25, measured 1.26–1.32;
width 64 takes 0.355); the affine per-MMA extrapolation under-predicted width 256 by 16–32 percent and is retired (measured,
MM~\mm{25.124.2}, MM~\mm{25.126}).

Simdgroups per threadgroup 1, 2, 4 and 8 predict the same attention time; the model keeps the measured default 8
(1.455~ms, 4.2 percent behind the best, 2 at 1.396) and excludes 16 (16 threadgroups leave 4 cores idle;
MM~\mm{25.124}). The model chooses schedules but does not lower them, and the production runtime cannot express most
schedules it chooses between (module docstring).

\subsection*{Static model from compiled bytes (\code{tools/g17perfmodel.py})}

The model simulates one simdgroup in order over decoded bytes (MM~\mm{25.144.12} for every figure here). Each
instruction issues one 3.105~ns interval after the last and not before its operands. Loads publish scoreboard slots
and consumers stall until they land (tensor fragment loads, \op{12674} and \op{12709}, take the working-set latency;
\op{12682} the 16 KiB latency), and the result of a 16-bit MMA, \op{5106}, is ready after a dependent-MMA latency.
Loops are walked 4 trips and extrapolated (latches \code{cnt < imm} and \code{p > 0} only). Threadgroups are dealt
to 20 cores in index order, and residency comes from the program's highest register (the MM~\mm{25.121} table) and
3,072 threads per core. Per-core rate is an exact mean-value analysis of three stations (issue slots, load path,
tensor unit) fed into the MM~\mm{25.124.1} list schedule. Memory is floored at 280~GB/s of DRAM, and each dispatch
adds a fixed cost per timing protocol.

The constants are fitted by bounded least squares on log error, with the tensor unit not allowed to beat 4.95~ns per
MMA per core: 12.5 issue slots per interval per core; 3.43~ns per fragment load (about 150~GB/s of fragments per
core); a scalar load at 0.43 fragment loads; 6.83~ns per MMA (72 percent of peak); load latency 0.50 times the
pointer chase; dependent MMA latency 31.6~ns; MMA issue
1.0 interval; per-dispatch cost 1.4 / 1.0 / 1.0~µs. The latency and issue constants sit at their lower bounds, so
these timings cannot identify them.

Over 108 recorded timings (45 GEMM, 54 attention, 9 graph; 16 set aside as noisy) the leave-one-table-out median error
is 14.8 percent against a ±15 percent target, with 47 of 92 within ±15 percent. A fixed split gives 15.5 percent and
the in-sample fit 9.6 percent. By subset the error is 8.3 percent for graph attention, 9.9 percent for hot attention
at M ≥ 512, 19.6 percent for warm GEMM and 35 percent for the noisy rows.

GEMM tile shapes do not extrapolate: a 2×2 tile is predicted up to 40 percent slow and tall tiles 16–37 percent fast,
because operand sharing among a threadgroup's simdgroups is invisible to the model. 1,024-token prefill projections
were predicted at 144~ms against 113 measured (+27 percent), and SwiGLU at 9.0 against 7.2~ms (DRAM-bound, and the
model has one 280~GB/s plateau). The model refuses to run without the
opcode class table, because an UNKNOWN fallback moved 81 of 108 predictions. Two fetch terms built from the recorded
cliff (1.91~ns per byte cold, 2.81 times past 16 KiB on one simdgroup, 1.15 times at 64) both worsened held-out error
(23.2 and 17.5 percent against 14.8), so fetch is not in the model of record.

The model predicted savings of 7.6~ms for SwiGLU fused into w3's GEMM and 1.5~ms for hardware exp2 at 512-row
chunks (679 to 490 instructions a trip), and both were then built. It predicted a 2.8~ms loss for sharing K and V
across 4 simdgroups through threadgroup memory.

\textbf{Evidence.} MM~\mm{16}, MM~\mm{25.124}, MM~\mm{25.124.1}, MM~\mm{25.124.2}, MM~\mm{25.126}, MM~\mm{25.144.12}; \code{agxforge/g17/tensorsched.py}.

\section{Model workloads}\label{sec:bounds-real-workloads}

The workload is InternLM2.5-1.8B at 4 bits unless stated. Throughput against mlx-lm for each workload is in the
machine model (MM~\mm{25.190}; the headline comparison is in \href{https://github.com/sbryngelson/AGXForge/blob/main/docs/demonstrations.md}{the demonstrations document}, section 2).

\subsection*{Single-stream decode}

In the MM~\mm{25.124.5} decode plan (model and measured), one bf16 decode step (d 2048, 16 heads × 128, KV 256, one
query row) has a memory floor of 0.487~ms (136.3~MB), 0.247~ms with fp8 weights. A one-token projection with too few
threadgroups is bound by its K chain (measured, MM~\mm{25.124.6}). M 16, N 128, K 2048 takes the same time with cold
and cache-resident weights (within 1 percent), at 6.5–9.9~GB/s for this compiler and 7.5–16.7~GB/s for Apple's,
against 270. With enough threadgroups the column grid reaches the floor. N 8192 × K 2048 (33.5~MB) on 64 threadgroups
of 128 columns takes 143–145~µs = 231–234~GB/s, 1.2 times its 0.120~ms floor, equal to Apple's \code{matmul2d}
(142.8–145.2~µs). Split-K on N 2048, K 2048 falls from 84.2~µs at G 1 to 52.1 at G 2 and 41.1 at G 4, then flattens
(42.0 at G 8, 43.3 at G 16); the plateau, about 200~GB/s with cached and cold weights alike, is therefore not set by
DRAM.

At a 1K context a token takes 4.73~ms, split by part in \cref{tab:decode-time-breakdown} (measured, MM~\mm{25.181}).

\begin{xltabular}{\linewidth}{@{}>{\hsize=1.20\hsize}Xl>{\hsize=0.80\hsize}Xl@{}}
\caption{Single-stream decode time per token at a 1K context.}\label{tab:decode-time-breakdown}\\
\toprule
Part & ms & Bytes per token & Effective GB/s \\
\midrule
\endfirsthead
\tablecontinued{4}
\toprule
Part & ms & Bytes per token & Effective GB/s \\
\midrule
\endhead
FFN (w1, w3, SwiGLU) & 1.76 & 453~MB & 257 \\*
w2 + residual & 0.90 & 227~MB & 252 \\*
qkv & 0.48 & 113~MB & 236 \\*
wo + residual & 0.30 & 57~MB & 188 \\*
head & 0.36 & 107~MB & 296 \\*
attention & 0.74 & 107~MB of KV & 144 \\*
norms and the rest & 0.18 &  &  \\*
\bottomrule
\end{xltabular}

Only the head streams at the DRAM peak. Counters on the decode qmv kernels (measured, MM~\mm{25.144.4}) show the q4 w2
kernel issue-bound (instruction throughput limiter 77.2 percent, zero memory-stall influence; pinned-weight floor
32.9~µs above the 29.8~µs DRAM floor) and q8 w2 memory-bound (instruction limiter 46.7 percent, memory-stall influence
26.6 percent, LLC limiter 100; 89 percent of the 298.5~GB/s peak). MLX's qmv is latency- and memory-bound at 16.7
percent occupancy (64-thread threadgroups), with 37–59 percent memory-stall influence.

In q4 w2, compute and stream barely overlap (measured, MM~\mm{25.141.15}, at the 4-rows-per-simdgroup shape): real
46.2~µs; x index cache-resident 40.9; weight index resident 29.5; both 28.5~µs (the instruction floor); streaming
8.9~MB at 298.5~GB/s needs about 30~µs. The single-stream q4 MLP is load-bound
(measured, MM~\mm{25.200}): removing 26 percent of its instructions (FFN 393 to 291 instructions, 10 loads either
way) does not change throughput (186.4–192.1 against 177.9–193.4 tok/s at 1,792 tokens).

Context costs about 20~ns per key-layer (measured, MM~\mm{25.199}), and long-context attention variants
(\code{keyblock}, \code{tgsplit}) do not help in the graph at 1,792 keys. Apple-compiled twins of this project's decode
kernels, bit-identical in output, give 10–15 percent higher median throughput, so this project's instruction-level
code generation costs that much (measured, MM~\mm{25.211}; demonstrations document, section 2).

\subsection*{Batched decode}

The tensor-unit projection (\code{g17qsm}: $Y^{T} = W^{T} X^{T}$, W dequantized once into registers as A, 16 sequences
as B) costs 446~µs per layer for 16 tokens (27.9~µs a token) against batched qmv's 126 at B~8, 4.5 times per token.
It runs at about 76~GB/s, latency-bound per simdgroup (about 700 dependent scalar instructions per trip before 8 MMAs;
MM~\mm{25.166}). Split-K to about 512 threadgroups is the knee at every role (qkv split 4: 39.8~µs at 512
threadgroups; wo split 8: 23.2; w1/w3 split 2: 76.4; w2 split 8: 75.8; MM~\mm{25.185}).

At B~16 the projections are bound by the fp32 pipe (F32 Limiter 100, F32 Utilization 99.9, instruction throughput
limiter 78 (p90 98), kernel occupancy 38 (p90 50), memory request stalls 0, GPU read bandwidth p90 122~GB/s, Neural
Accelerator utilization p90 25, F16 utilization 0). Projections are 7.7~ms of a 10.98~ms step (70 percent) for
0.95~GB of q4, about 2.6 times the DRAM floor. An ablation on w1 gives base 76.3, no dequant 59.6, no weight load 60.9
and neither 35.1~µs (MM~\mm{25.194}). Moving the dequant to the fp16 pipe with \op{798}, two per weight, relieves the
fp32 pipe, taking w1/w3 from 74.6 to 65.7~µs (−12 percent), w2 from 73.9 to 66.6, qkv from 39.6 to 35.4 and wo from
22.1 to 20.0 (MM~\mm{25.196}).

At B~8 the w2 projection hits the LLC limiter (100 percent) for this compiler and for MLX's \code{affine\_qmv\_wide}
alike, and the B~8 qkv is register-capped at 25–27 percent occupancy (\cref{sec:registers-spills-occupancy};
MM~\mm{25.144.4}). MLX's own batched decode is compute-bound past B~2 (quantized) and B~8 (bf16); every precision
converges on 19–21~ms per step at B~16 (q4 841 tok/s, bf16 750), and dequantization costs more than the bytes it
saves (MM~\mm{25.142.5}). The fastest mlx-lm batched path depends on the model. For InternLM2 it is the bench loop,
since \code{BatchGenerator} cannot run that model, and for Qwen3 \code{BatchGenerator} is 7–9 percent faster at B~8
and 16 (MM~\mm{25.193}).

\subsection*{Prefill}

With a closure check (dispatch-alone sum 135.8~ms against chained step spans 137.1 and in-process medians 139.6,
within 1 percent), the in-graph prefill GEMMs are at 80–89 percent of 33.1 TFLOP/s: qkv 256 × 4096 × 2048 148~µs (88
percent), wo 81~µs (80), w1 290~µs (89), w2 292~µs (89). The earlier figure of 60–69 percent, against MLX's 82–88,
came from a per-dispatch profile that did not close (181~ms of GPU time in a 150~ms prefill) and an isolated GEMM
bench that runs about four times the in-graph per-call time (\cref{app:a}). Pinning the weights cache-resident changed
per-role time by 0 to ±5 percent, so weight traffic is not the GEMM cost (MM~\mm{25.160}).

Prefill attention is latency-bound on the per-block scratch round trip: 32-key blocks cut instructions per 16 keys
from 490 to about 390 (−20 percent) for −3 percent time (MM~\mm{25.162}). Keeping S and P in registers cut attention
from 14.7 to 12.1~ms with exact exp, or to 9.6~ms with hardware exp2 (MM~\mm{25.163}). The SwiGLU kernel was bandwidth-bound, since 10
bytes an element over 201 M elements is 6.7~ms at 298.5~GB/s, and it measured 6.74~ms; an fp16 intermediate (6 bytes an
element) made it 5.54~ms and arithmetic-bound (MM~\mm{25.183}).

\begin{xltabular}{\linewidth}{@{}X>{\hsize=1.34\hsize}X>{\hsize=0.69\hsize}X@{}}
\caption{Binding resource for each studied workload.}\label{tab:binding-resources}\\
\toprule
Workload & Binding resource in the studied kernels & Evidence \\
\midrule
\endfirsthead
\tablecontinued{3}
\toprule
Workload & Binding resource in the studied kernels & Evidence \\
\midrule
\endhead
fused attention, FFN layer kernels & instruction issue and ALU (about 46 instructions per MMA) & MM~\mm{8}, MM~\mm{16} \\*
single-stream decode, q8 projections and the head & DRAM bandwidth & MM~\mm{25.144.4}, MM~\mm{25.181} \\*
single-stream decode, q4 w2 & instruction issue by counter, load-bound by ablation & MM~\mm{25.144.4}, MM~\mm{25.200} \\*
one-token projection with few threadgroups & one K chain's latency & MM~\mm{25.124.6} \\*
decode attention past 4.2~MB of KV & LLC capacity, then about 140~GB/s & MM~\mm{25.181} \\*
batched decode B~16 projections & occupancy first (about 256 one-simdgroup threadgroups a projection);
after split-K to about 512, the fp32 pipe, relieved by moving the
dequant to the fp16 pipe & MM~\mm{25.185}, MM~\mm{25.194}, MM~\mm{25.196} \\*
batched decode B~8 w2 & LLC bandwidth (both stacks) & MM~\mm{25.144.4} \\*
prefill GEMMs & tensor and issue, at 80–89 percent of peak & MM~\mm{25.160} \\*
prefill attention & latency of a scratch round trip, until S and P stay in registers & MM~\mm{25.162}, MM~\mm{25.163} \\*
\bottomrule
\end{xltabular}

\textbf{Evidence.} MM~\mm{8}, MM~\mm{16}, MM~\mm{25.124.5}, MM~\mm{25.124.6}, MM~\mm{25.141.15}, MM~\mm{25.142.5}, MM~\mm{25.144.4}, MM~\mm{25.160}, MM~\mm{25.162},
MM~\mm{25.163}, MM~\mm{25.166}, MM~\mm{25.181}, MM~\mm{25.183}, MM~\mm{25.185}, MM~\mm{25.190}, MM~\mm{25.193}, MM~\mm{25.194}, MM~\mm{25.196}, MM~\mm{25.199},
MM~\mm{25.200}, MM~\mm{25.211}.
